% =====================================================================
%  Thème 2 — Chasles et combinaisons — Niveau MOYEN — Exercices
% =====================================================================
\documentclass[11pt,a4paper]{article}
\usepackage[utf8]{inputenc}
\usepackage[T1]{fontenc}
\usepackage{lmodern}
\usepackage[french]{babel}
\usepackage{amsmath,amssymb,mathtools}
\usepackage{xcolor}
\usepackage{enumitem}
\usepackage{fancyhdr}
\usepackage[a4paper,margin=2.1cm,headheight=26pt,headsep=10pt]{geometry}

\definecolor{primary}{RGB}{222,70,80}
\definecolor{primarydark}{RGB}{150,30,42}

\pagestyle{fancy}\fancyhf{}
\fancyhead[L]{\small\textcolor{primarydark}{\textbf{Fiche 2} — Calcul vectoriel}}
\fancyhead[R]{\small\textcolor{primarydark}{\textsc{Thème 2 — Moyen — Exercices}}}
\fancyfoot[C]{\small\thepage}
\renewcommand{\headrulewidth}{0.8pt}

\newcommand{\vc}[1]{\overrightarrow{#1}}
\providecommand{\norm}[1]{\left\lVert #1\right\rVert}
\newcommand{\exo}[1]{\medskip\noindent{\color{primary}\bfseries\large Exercice #1.}\par\nopagebreak\smallskip}

\begin{document}
\begin{center}
{\LARGE\bfseries\color{primarydark}Chasles et combinaisons — Niveau Moyen}\\[2pt]
{\color{primary}\rule{0.5\linewidth}{1.1pt}}
\end{center}
\vspace{1mm}
{\small\itshape Objectif : réordonner et transformer avant d'appliquer Chasles ; introduire un point,
manipuler des milieux et des combinaisons linéaires.}
\vspace{2mm}

\exo{1} \textbf{Simplifier en un seul vecteur.}
\begin{enumerate}[label=\textbf{\alph*)},leftmargin=2em,itemsep=2pt]
  \item $\vc{AB}+\vc{CD}+\vc{BC}$
  \item $\vc{AB}-\vc{AC}$
  \item $\vc{AB}+\vc{BC}-\vc{DC}$
  \item $\vc{OA}-\vc{OB}$
\end{enumerate}

\exo{2} \textbf{Introduire un point $O$.}
Soit $O$ un point quelconque du plan. Exprimer les vecteurs suivants uniquement à l'aide de
$\vc{OA}$, $\vc{OB}$ et $\vc{OC}$.
\begin{enumerate}[label=\textbf{\alph*)},leftmargin=2em,itemsep=2pt]
  \item $\vc{AB}$
  \item $\vc{CA}$
  \item $\vc{AB}+\vc{BC}$
\end{enumerate}

\exo{3} \textbf{Milieu.}
$I$ est le milieu du segment $[AB]$.
\begin{enumerate}[label=\textbf{\alph*)},leftmargin=2em,itemsep=2pt]
  \item Exprimer $\vc{AI}$ en fonction de $\vc{AB}$.
  \item Que vaut $\vc{IA}+\vc{IB}$ ? Justifier.
  \item Montrer que, pour tout point $O$, $\vc{OI}=\dfrac12\bigl(\vc{OA}+\vc{OB}\bigr)$.
\end{enumerate}

\exo{4} \textbf{Diagonales d'un parallélogramme.}
$ABCD$ est un parallélogramme. On pose $\vec u=\vc{AB}$ et $\vec v=\vc{AD}$.
\begin{enumerate}[label=\textbf{\alph*)},leftmargin=2em,itemsep=2pt]
  \item Exprimer la diagonale $\vc{AC}$ en fonction de $\vec u$ et $\vec v$.
  \item Exprimer la diagonale $\vc{BD}$ en fonction de $\vec u$ et $\vec v$.
\end{enumerate}

\exo{5} \textbf{Combinaison linéaire avec un milieu.}
$ABCD$ est un parallélogramme et $I$ est le milieu du côté $[BC]$. On pose $\vec u=\vc{AB}$ et
$\vec v=\vc{AD}$.
\begin{enumerate}[label=\textbf{\alph*)},leftmargin=2em,itemsep=2pt]
  \item Exprimer $\vc{BI}$ en fonction de $\vec v$.
  \item En déduire $\vc{AI}$ comme combinaison linéaire de $\vec u$ et $\vec v$.
\end{enumerate}

\end{document}
